\textbf{1. Kiến thức cần nhớ}
\begin{itemize}
\item \textbf{Hình chóp tam giác đều} .ABC$ là hình chóp có:
\begin{itemize}
\item Đáy $ là một \textbf{tam giác đều}.
\item Các mặt bên , SBC, SCA$ là các \textbf{tam giác cân bằng nhau} tại đỉnh $.
\item Các cạnh bên bằng nhau:  = SB = SC$. Các cạnh đáy bằng nhau:  = BC = CA$.
\end{itemize}
\item \textbf{Trung đoạn} của hình chóp tam giác đều là đoạn thẳng nối đỉnh $ với trung điểm của một cạnh đáy (chính là đường cao của tam giác mặt bên kẻ từ đỉnh $, kí hiệu là $).
\item \textbf{Đường cao} là đoạn thẳng nối đỉnh $ với trọng tâm $ của tam giác đáy ( \perp (ABC)$, độ dài  = h$).
\item \textbf{Diện tích xung quanh}: Bằng nửa tích chu vi đáy với trung đoạn (hoặc tổng diện tích $ mặt bên):
80830\boxed{S_{xq} = p \cdot d = 3 \cdot \left(\dfrac{1}{2} a \cdot d\right)}80830
(với $ là nửa chu vi đáy, $ là độ dài cạnh đáy, $ là độ dài trung đoạn).
\item \textbf{Diện tích toàn phần}: Bằng tổng diện tích xung quanh và diện tích đáy:
80830\boxed{S_{tp} = S_{xq} + S_{\text{đáy}}}80830
\item \textbf{Thể tích}: Bằng một phần ba tích diện tích đáy với chiều cao của hình chóp:
80830\boxed{V = \dfrac{1}{3} S_{\text{đáy}} \cdot h}80830
(với {\text{đáy}} = \dfrac{a^2\sqrt{3}}{4}$ là diện tích tam giác đều cạnh $,  = SO$ là chiều cao).
\end{itemize}

\textbf{2. Phương pháp giải các dạng}

%% DANG 8H0?38-1
\textbf{Dạng 1. Nhận biết các yếu tố của hình chóp tam giác đều}
\begin{enumerate}
\item Chỉ ra đỉnh, mặt đáy (tam giác đều), các mặt bên (các tam giác cân bằng nhau).
\item Đếm số đỉnh ($ đỉnh), số mặt ($ mặt: $ đáy và $ mặt bên), số cạnh ($ cạnh: $ cạnh đáy và $ cạnh bên).
\item Nhận biết hình khai triển trên mặt phẳng (gồm $ tam giác đều ở giữa và $ tam giác cân chung cạnh với tam giác đều).
\end{enumerate}

%% DANG 8H0?38-2
\textbf{Dạng 2. Diện tích xung quanh, thể tích hình chóp tam giác đều}
\begin{enumerate}
\item Xác định các đại lượng: cạnh đáy $, trung đoạn $, chiều cao $.
\item Sử dụng định lí Pythagore trong tam giác vuông để tìm mối liên hệ giữa các đại lượng:
\begin{itemize}
\item Trong tam giác mặt bên: ^2 + \left(\dfrac{a}{2}\right)^2 = \text{(cạnh bên)}^2$.
\item Trong tam giác vuông chứa đường cao $: ^2 + R^2 = \text{(cạnh bên)}^2$ (với  = \dfrac{a\sqrt{3}}{3}$ là bán kính đáy).
\end{itemize}
\item Thay số vào công thức {xq} = \dfrac{3}{2}a \cdot d$ và  = \dfrac{1}{3}S_{\text{đáy}} \cdot h$.
\end{enumerate}

%% DANG 8H0?38-3
\textbf{Dạng 3. Bài toán thực tế}
\begin{enumerate}
\item Mô hình hóa các vật thể trong đời sống (hộp quà, lều cắm trại, đèn kéo quân, mái nhà\dots) về hình chóp tam giác đều.
\item Xác định yêu cầu đề bài:
\begin{itemize}
\item Làm bằng bạt/vải/giấy xung quanh (không tính đáy) $\rightarrow$ tính {xq}$.
\item Sơn/dán toàn bộ cả đáy $\rightarrow$ tính {tp}$.
\item Chứa lượng không khí/dung tích bên trong $\rightarrow$ tính $.
\end{itemize}
\item Tính toán chi phí, lượng vật liệu kèm hao hụt (nếu đề bài cho).
\end{enumerate}
